\newpage
\section{Conclusion et perspectives}
\label{sec:conclusion}

\subsection{Résumé du travail réalisé}
\paragraph{} Le projet "Tu n'y peux rien" a été **mené à bien dans les délais impartis et a été entièrement achevé**. Toutes les fonctionnalités demandées dans le cahier des charges ont été implémentées et testées :

\begin{itemize}
\item Jeu complet avec interface graphique Swing fluide et responsive
\item 3 niveaux de difficulté pour les bots (Facile, Moyen, Difficile)
\item Détection robuste des combinaisons (simple, double, série, bombe, double joker)
\item Gestion correcte des cartes spéciales (2, Joker)
\item Système de scores et statistiques avec graphiques jFreeChart
\item Tests unitaires (JUnit) couvrant les classes critiques
\item Logs détaillés (Log4j) en HTML et texte
\end{itemize}

\paragraph{} Au total, ce sont **plus de 40 classes Java** qui ont été développées, pour environ **5000 lignes de code**. Le projet compile et s'exécute sans erreur sur les environnements cibles.

\subsection{Tests et validation}
\label{sec:tests}

\paragraph{} Les tests unitaires ont été réalisés avec JUnit 4. Chaque classe critique possède sa classe de test associée :

\begin{itemize}
\item \texttt{TestCard.java} : test de la création des cartes, des symboles et de l'égalité
\item \texttt{TestDeck.java} : test du mélange et de la pioche
\item \texttt{TestPlayer.java} : test de la pioche et de la pose de combinaisons
\item \texttt{TestCombination.java} : test de la détection des types et de la comparaison (15 cas)
\item \texttt{TestGameStats.java} : test du système de scores
\end{itemize}

\paragraph{} Au total, \textbf{37 tests unitaires} ont été écrits et passent tous avec succès. La couverture de code est satisfaisante sur les classes métier (\texttt{Combination}, \texttt{Player}, \texttt{GameStats}) mais reste partielle sur l'interface graphique (tests manuels uniquement).

\subsection{Améliorations possibles du projet}
\paragraph{} Bien que le projet soit fonctionnel et complet, plusieurs extensions pourraient enrichir l'expérience utilisateur :

\begin{itemize}
\item \textbf{Mode multijoueur en ligne} : Permettre à plusieurs joueurs humains de jouer à distance via réseau (sockets Java)
\item \textbf{Sauvegarde des parties} : Sauvegarder l'état du jeu dans un fichier pour le reprendre plus tard
\item \textbf{Classements et statistiques persistants} : Stocker les scores et victoires dans une base de données SQLite
\item \textbf{Animations} : Ajouter des animations lors de la pose des cartes (transition, mouvement)
\item \textbf{Thèmes personnalisables} : Permettre au joueur de changer l'apparence des cartes et du tapis (plusieurs skins)
\item \textbf{Version mobile} : Adapter le jeu pour Android via un wrapper ou une réécriture partielle
\end{itemize}

\paragraph{} En conclusion, le projet "Tu n'y peux rien" remplit l'ensemble des objectifs fixés. Il démontre la maîtrise des concepts de génie logiciel (Java, Swing, patterns, tests, Git) et offre une expérience de jeu complète et agréable. Ces améliorations constitueraient une excellente base pour une version 2.0 du projet.

\newpage
\section{Annexes}
\label{sec:annexes}

\subsection{Guide d'installation rapide}
\paragraph{Prérequis :}
\begin{itemize}
\item Java 8 ou supérieur (JDK 1.8+)
\item Eclipse IDE (version 2020-06 ou supérieure)
\end{itemize}

\paragraph{Import du projet dans Eclipse :}
\begin{enumerate}
\item Ouvrir Eclipse
\item \texttt{File $\rightarrow$ Import $\rightarrow$ Git $\rightarrow$ Projects from Git}
\item Sélectionner \texttt{Clone URI}
\item Renseigner : \texttt{https://github.com/berrandou-dev/Projet\_GLP-Jeu\_Carte.git}
\item Suivre les instructions par défaut
\item Une fois importé, faire \texttt{Project $\rightarrow$ Clean} pour forcer la compilation
\end{enumerate}

\paragraph{Ajout des bibliothèques externes :}
\begin{enumerate}
\item Faire clic droit sur le projet $\rightarrow$ \texttt{Build Path $\rightarrow$ Configure Build Path}
\item Onglet \texttt{Libraries} $\rightarrow$ \texttt{Add JARs...}
\item Sélectionner tous les JARs dans \texttt{src/lib/} :
\begin{verbatim}
- log4j-1.2.17.jar
- junit-4.11.jar
- jfreechart-1.0.19.jar
- jcommon-1.0.23.jar
\end{verbatim}
\item Cliquer sur \texttt{Apply and Close}
\end{enumerate}

\paragraph{Lancement du jeu :}
\begin{enumerate}
\item Dans l'Explorateur de projet, naviguer vers \texttt{src/test/manual/}
\item Faire clic droit sur \texttt{TestGame.java} $\rightarrow$ \texttt{Run As $\rightarrow$ Java Application}
\item Le menu principal apparaît
\end{enumerate}

\subsection{Structure des packages du projet}
\paragraph{} Le projet est organisé selon l'architecture MVC :

\begin{verbatim}
src/
+-- config/           # Configuration (GameConfig.java)
+-- engine/
    +-- data/         # MODELE : Card, Deck, Player, GameStats
    +-- process/      # CONTROLEUR : Game, Combination, bots
+-- gui/
    +-- panels/       # VUE : composants reutilisables
    +-- screens/      # VUE : ecrans complets
    +-- utils/        # Utilitaires : GameStyle, MusicPlayer, ChartManager
+-- lib/              # Bibliotheques externes (JARs)
+-- log/              # Logs generes (ignores par Git)
+-- resources/        # Images et musique
+-- test/             # Tests unitaires (JUnit)
    +-- demo/         # Mode demonstration
    +-- manual/       # Tests manuels (TestGame, TestEndGameGUI)
    +-- unit/         # Tests unitaires par classe
\end{verbatim}

\subsection{Principales classes du projet}
\paragraph{} Voici un résumé des classes les plus importantes :

\begin{table}[h!]
\centering
\begin{tabular}{|l|l|p{6cm}|}
\hline
\textbf{Classe} & \textbf{Package} & \textbf{Rôle} \\
\hline
\texttt{Card} & engine.data & Représente une carte (valeur + couleur) \\
\texttt{Deck} & engine.data & Gère le paquet de 54 cartes \\
\texttt{Player} & engine.data & Gère la main du joueur \\
\texttt{GameStats} & engine.data & Enregistre les scores et statistiques \\
\hline
\texttt{Game} & engine.process & Logique principale du jeu \\
\texttt{Combination} & engine.process & Détection et comparaison des combinaisons \\
\texttt{BotPlayer} & engine.process & Classe abstraite pour les bots \\
\hline
\texttt{MainGUI} & gui.screens & Interface principale du jeu \\
\texttt{EndGameGUI} & gui.screens & Écran de fin de partie avec graphiques \\
\texttt{ChartManager} & gui.utils & Génération des graphiques jFreeChart \\
\hline
\end{tabular}
\caption{Principales classes du projet}
\label{tab:classes}
\end{table}
